% =====================================================================
%  Intro to LaTeX: Formel-Challenge (Aufgabe 2 aus dem Praxisteil)
%  jDPG Regionalgruppe Tübingen
%  Seite 1: Aufgaben, Seite 2: Lösungen
%  Kompilieren mit pdfLaTeX: latexmk -pdf formel-challenge.tex
% =====================================================================
\documentclass[a4paper,11pt]{scrartcl}

% ---------------------------------------------------------------------
%  Pakete
% ---------------------------------------------------------------------
\usepackage[T1]{fontenc}
\usepackage[ngerman]{babel}
\IfFileExists{FiraSans.sty}{%
  \usepackage[sfdefault,lining]{FiraSans}
  \usepackage[scale=0.85]{FiraMono}}{%
  \IfFileExists{lmodern.sty}{\usepackage{lmodern}}{}}
\usepackage[a4paper, top=2cm, bottom=2cm, left=2.2cm, right=2.2cm]{geometry}
\usepackage{amsmath,amssymb}
\usepackage{listings}
\usepackage[most]{tcolorbox}

\setlength{\parindent}{0pt}
\setlength{\parskip}{0.4em}
\pagestyle{empty}

% ---------------------------------------------------------------------
%  Farben (wie in den Folien)
% ---------------------------------------------------------------------
\definecolor{dunkel}{HTML}{16213E}
\definecolor{akzent}{HTML}{E9A23B}
\definecolor{blau}{HTML}{3D7EA6}
\definecolor{hell}{HTML}{EDF1F5}
\definecolor{grau}{HTML}{6B7280}
\definecolor{codebg}{HTML}{F2F4F7}

% ---------------------------------------------------------------------
%  Bausteine
% ---------------------------------------------------------------------
% Titelzeile wie im Handout
\newcommand{\titelzeile}[1]{%
  \begin{tcolorbox}[enhanced, colback=dunkel, colframe=dunkel, arc=2pt,
    boxrule=0pt, left=10pt, right=10pt, top=7pt, bottom=7pt, before skip=0pt]
    \color{white}
    {\sffamily\bfseries\fontsize{20}{22}\selectfont Formel-Challenge}\hfill
    {\sffamily\small\color{white!70!dunkel} #1}\par
    \vspace{2pt}{\color{akzent}\rule{2.2cm}{2pt}}\hfill
    {\scriptsize\color{white!60!dunkel} Intro to \LaTeX{} \,|\, jDPG Regionalgruppe Tübingen}
  \end{tcolorbox}}

% Level-Markierung wie \pill in den Folien
\newcommand{\level}[2]{%
  \par\vspace{1em}%
  \tikz[baseline=(p.base)]\node[fill=#1, text=white, rounded corners=6pt,
    inner xsep=6pt, inner ysep=2.5pt, font=\small\bfseries] (p) {Level #2};\par}

\newtcolorbox{formel}{enhanced, colback=white, colframe=hell, boxrule=0.8pt,
  arc=4pt, halign=center, top=6pt, bottom=6pt}

\lstdefinestyle{latex}{
  language=[LaTeX]TeX,
  basicstyle=\ttfamily\small\color{dunkel},
  commentstyle=\color{grau}\itshape,
  texcsstyle=*\color{blau}\bfseries,
  columns=fullflexible, keepspaces=true, breaklines=true,
  showstringspaces=false}
\newtcblisting{code}{enhanced, listing only, colback=codebg, colframe=codebg,
  boxrule=0pt, arc=4pt, left=8pt, right=6pt, top=2pt, bottom=2pt,
  listing options={style=latex}}

% Formeln nur einmal definieren, damit Aufgabe und Lösung gleich bleiben
\newcommand{\formelA}{F = m a, \qquad E = m c^2, \qquad p = \frac{h}{\lambda}}
\newcommand{\formelB}{\nabla \times \vec{B} = \mu_0 \vec{j} + \mu_0 \varepsilon_0 \frac{\partial \vec{E}}{\partial t}}
\newcommand{\formelC}{\langle x \rangle = \int_{-\infty}^{\infty} \psi^*(x)\, x\, \psi(x) \,\mathrm{d}x}
\newcommand{\formelD}{Z = \sum_{n=0}^{\infty} \exp\left(-\frac{E_n}{k_\mathrm{B} T}\right), \qquad E_n = \hbar\omega\left(n + \frac{1}{2}\right)}

\begin{document}

% =====================================================================
%  Aufgaben
% =====================================================================
\titelzeile{Aufgaben}

Setzt die Formeln so genau wie möglich nach. Achtet auf Kleinigkeiten:
aufrechte Indizes, Abstände, Klammergrößen.

\level{blau!60}{1}
\begin{formel} \vspace{-1.2em}\[ \formelA \] \end{formel}

\level{blau}{2}
\begin{formel} \vspace{-1.2em}\[ \formelB \] \end{formel}

\level{akzent}{3}
\begin{formel} \vspace{-1.2em}\[ \formelC \] \end{formel}

\level{akzent!70!red}{4}
\begin{formel} \vspace{-1.2em}\[ \formelD \] \end{formel}

\vfill
{\small\color{grau} Tipp: Symbol unbekannt? Auf \texttt{detexify.kirelabs.org}
zeichnen und den Befehl bekommen. Die Lösungen stehen auf der nächsten Seite.}

\newpage
% =====================================================================
%  Lösungen
% =====================================================================
\titelzeile{Lösungen}

\level{blau!60}{1}
\begin{code}
\[ F = m a, \qquad E = m c^2, \qquad p = \frac{h}{\lambda} \]
\end{code}
\textbf{Achtung:} \texttt{\textbackslash qquad} für Abstand zwischen Formeln,
nicht viele Leerzeichen (die ignoriert \LaTeX{} im Mathemodus).

\level{blau}{2}
\begin{code}
\[ \nabla \times \vec{B} = \mu_0 \vec{j}
   + \mu_0 \varepsilon_0 \frac{\partial \vec{E}}{\partial t} \]
\end{code}
\textbf{Achtung:} \texttt{\textbackslash varepsilon} ($\varepsilon$) statt
\texttt{\textbackslash epsilon} ($\epsilon$), \texttt{\textbackslash times} für das Kreuzprodukt.

\level{akzent}{3}
\begin{code}
\[ \langle x \rangle = \int_{-\infty}^{\infty}
   \psi^*(x)\, x\, \psi(x) \,\mathrm{d}x \]
\end{code}
\textbf{Achtung:} \texttt{\textbackslash langle}, \texttt{\textbackslash rangle} statt
\texttt{<}, \texttt{>}. Das Differential $\mathrm{d}x$ aufrecht und mit kleinem
Abstand \texttt{\textbackslash,} davor.

\level{akzent!70!red}{4}
\begin{code}
\[ Z = \sum_{n=0}^{\infty} \exp\left(-\frac{E_n}{k_\mathrm{B} T}\right),
   \qquad E_n = \hbar\omega\left(n + \frac{1}{2}\right) \]
\end{code}
\textbf{Achtung:} \texttt{\textbackslash left(} \dots\ \texttt{\textbackslash right)}
passt die Klammern an den Bruch an. Der Index B in $k_\mathrm{B}$ ist ein Name
(Boltzmann), keine Variable, deshalb aufrecht mit \texttt{\textbackslash mathrm}.

\end{document}
